\section{Síntesis y ampliación}

Esta clase replantea el security leadership: deja de ser la pregunta «¿es el CSO lo bastante valiente?» y se convierte en un diseño organizacional verificable:

\begin{enumerate}
  \item \textbf{La seguridad es un operating system}: people, process, platform y governance determinan juntos el resultado;
  \item \textbf{Las reglas públicas necesitan retroalimentación}: las empresas deben llevar la evidence de sistemas reales a la elaboración de políticas y, al mismo tiempo, aceptar la accountability pública;
  \item \textbf{El VDP gestiona el canal de reportes}: scope, acknowledgement, corrección y divulgación coordinada conforman un contrato de confianza;
  \item \textbf{El Bounty gestiona incentivos opcionales}: una recompensa, un NDA o un payment no pueden autorizar a posteriori una conducta dañina;
  \item \textbf{El Incident command gestiona la ejecución interfuncional}: el IC sincroniza los flujos de trabajo, pero no monopoliza los juicios técnicos, legales ni de negocio;
  \item \textbf{La Evidence y el decision log garantizan que todo pueda revisarse}: hechos, análisis, decisiones y correcciones deben registrarse en capas separadas;
  \item \textbf{Materiality no equivale a severity}: para las empresas cotizadas a las que aplica, el plazo de cuatro días hábiles del Form 8-K suele contarse a partir de que se determina la importancia del incidente;
  \item \textbf{El estado de los casos debe actualizarse}: la pending appeal mencionada en clase quedó superada por la sentencia del Noveno Circuito de 2025 y por la cert denial de 2026;
  \item \textbf{La Accountability se comparte, pero no se diluye}: CSO, GC, CEO y board tienen cada uno un deliverable y un decision right claramente definidos;
  \item \textbf{La Resilience y el public service son capacidades de largo plazo}: recuperar al personal, conservar la memoria institucional y mejorar la gobernanza pública de la tecnología.
\end{enumerate}

\begin{importantbox}{Tarea de Security Governance}
Diseña un playbook de extremo a extremo para una empresa SaaS ficticia: publica el VDP y los límites del safe-harbor; define los VDP-to-incident triggers; traza la matriz RACI de IC, technical, legal/comms y business; proporciona plantillas de evidence chain y de decision-log; enumera los workstreams de divulgación a customer, regulator, law-enforcement y public-market; por último, diseña seis tabletop injects. Cada decisión clave debe indicar owner, evidence, deadline y review condition.
\end{importantbox}

\begin{knowledgebox}{Autoevaluación final: cinco pares que no deben confundirse}
No confundas vulnerability con incident, ni severity con legal materiality, ni payment con authorization, ni speaker account con adjudicated outcome, ni shared responsibility con diffuse accountability. Mientras estos cinco límites estén claros, es más probable que la organización mantenga íntegros los hechos y las responsabilidades bajo presión.
\end{knowledgebox}

\subsection{Lecturas adicionales}

Los materiales siguientes no están en el mismo nivel: CISA, DOJ y NIST ayudan principalmente a diseñar el VDP y el incident-response process; los documentos de la SEC explican un marco específico de divulgación en mercados públicos; las opiniones judiciales y el docket registran el resultado procesal de los casos. Al leerlos, primero determina si quieres responder «cómo construir controles», «qué obligaciones aplican» o «qué juicio emitió el tribunal sobre el expediente existente»; así evitarás sustituir el análisis jurídico por guías de ingeniería, y también sustituir principios generales de ingeniería por un solo caso.

{\small
\begin{itemize}[nosep,leftmargin=*]
  \item \href{https://www.cisa.gov/vulnerability-disclosure-policy-template}{\textbf{CISA Vulnerability Disclosure Policy Template}}: ejemplos de scope, reporting, expectations y coordinated disclosure.
  \item \href{https://www.cisa.gov/news-events/directives/bod-20-01-develop-and-publish-vulnerability-disclosure-policy}{\textbf{CISA BOD 20-01}}: contexto de política del VDP para las agencias federales.
  \item \href{https://www.justice.gov/criminal/criminal-ccips/page/file/983996/dl}{\textbf{DOJ Framework for a Vulnerability Disclosure Program}}: marco sobre autorización, good faith y consideraciones de aplicación de la ley.
  \item \href{https://csrc.nist.gov/pubs/sp/800/61/r3/final}{\textbf{NIST SP 800-61 Rev. 3}}: integra la incident response en las actividades de gestión de riesgos del CSF 2.0.
  \item \href{https://www.sec.gov/files/33-11216-fact-sheet.pdf}{\textbf{SEC Cybersecurity Disclosure Final Rule Fact Sheet}}: panorama del material incident y de los plazos del Form 8-K.
  \item \href{https://cdn.ca9.uscourts.gov/datastore/opinions/2025/11/12/23-927.pdf}{\textbf{United States v. Sullivan, No. 23-927}}: opinión modificada del Noveno Circuito.
  \item \href{https://www.supremecourt.gov/docket/docketfiles/html/public/25-1082.html}{\textbf{Supreme Court Docket 25-1082}}: registro de la certiorari petition y de su rechazo del 29 de junio de 2026.
\end{itemize}
}

\begin{knowledgebox}{Orden de lectura sugerido: de los controles ejecutables a los límites judiciales}
Primero usa el CISA template para revisar si la policy pública establece scope, compromisos y canales de contacto; después usa el DOJ framework para examinar los límites de good-faith y de autorización; luego usa NIST SP 800-61 Rev. 3 para escalar el reporte a una incident response de nivel organizacional; si la organización es una empresa cotizada estadounidense a la que aplica la regla, lee el SEC fact sheet y el texto original de la final rule; por último, contrasta la Ninth Circuit opinion con el Supreme Court docket para confirmar el estado procesal del caso visto en clase. En cada paso registra la fecha de la fuente y su ámbito de aplicación, porque las políticas, las normas y el estado de los casos seguirán cambiando.
\end{knowledgebox}

\begin{importantbox}{Ejercicio de Source triangulation}
Elige un incident hipotético y escribe por separado: las recomendaciones operativas que el equipo técnico obtiene de NIST, las obligaciones de notificación y divulgación que counsel debe verificar, las decisiones de negocio que debe tomar la dirección y las preguntas que los materiales judiciales \textbf{no pueden} responder por ti. Si la misma conclusión aparece en las cuatro columnas, eso suele indicar que los niveles de evidencia siguen mezclados.
\end{importantbox}

\end{document}
